\section{业务流程分析}

\subsection{概述}

本节从参与者协作和业务对象状态变化两个角度描述系统的主要业务流程。外卖履约流程覆盖顾客下单、商家处理、骑手配送和管理员异常治理；智能点餐流程覆盖自然语言输入、受控查询、候选校验和用户确认。活动图说明跨角色工作流，状态机图说明对象生命周期，顺序图说明关键交互次序。

\subsection{外卖订单履约流程}

\subsubsection{总体活动流程}

\srsfigure{figures/srs/04-fulfillment-activity.pdf}{外卖订单履约活动图}{UML 活动图}{fig:fulfillment-activity}

图~\ref{fig:fulfillment-activity} 使用顾客、订单服务、商家、配送服务、骑手和通知服务六个泳道，包含库存不足、支付失败、商家拒单、用户取消和配送异常等分支。

\subsubsection{订单状态}

\srsfigure{figures/srs/05-order-state.pdf}{订单状态机图}{UML 状态机图}{fig:order-state}

\begin{longtable}{|L{0.16\linewidth}|L{0.18\linewidth}|L{0.18\linewidth}|L{0.38\linewidth}|}
\caption{订单状态转换}\label{tab:order-state}\\ \hline
\textbf{当前状态} & \textbf{事件/执行者} & \textbf{目标状态} & \textbf{条件与副作用} \\ \hline
无 & 提交订单/顾客 & 待支付 & 校验店铺、商品、地址和优惠；生成快照并预占库存与优惠券。 \\ \hline
待支付 & 支付成功/顾客 & 待商家接单 & 确认支付流水和资产核销，向商家发送新订单通知。 \\ \hline
待支付 & 主动取消或 15 分钟超时/顾客或系统 & 已取消 & 释放预占库存和优惠券，不生成扣款。 \\ \hline
待商家接单 & 接单/商家 & 制作中 & 校验订单属于本店且未取消，记录接单时间。 \\ \hline
待商家接单 & 拒单或顾客取消 & 已取消 & 记录取消方和原因，恢复库存并原路返还资产。 \\ \hline
制作中 & 确认出餐/商家 & 待骑手接单 & 生成一条待接配送任务；顾客不再可自助取消。 \\ \hline
待骑手接单 & 接单/骑手 & 待取餐 & 通过带状态条件的原子更新绑定唯一骑手。 \\ \hline
待取餐 & 确认取餐/骑手 & 配送中 & 仅本任务骑手可操作，记录取餐时间。 \\ \hline
配送中 & 确认送达/骑手 & 已送达 & 记录送达时间，向顾客发送确认收货提示。 \\ \hline
已送达 & 确认收货或 24 小时超时/顾客或系统 & 已完成 & 纳入营业额和销量统计，发放积分并开放评价入口。 \\ \hline
\end{longtable}

\subsubsection{配送任务状态}

\srsfigure{figures/srs/06-delivery-state.pdf}{配送任务状态机图}{UML 状态机图}{fig:delivery-state}

\begin{longtable}{|L{0.18\linewidth}|L{0.18\linewidth}|L{0.20\linewidth}|L{0.34\linewidth}|}
\caption{配送任务状态转换}\label{tab:delivery-state}\\ \hline
\textbf{当前状态} & \textbf{操作} & \textbf{目标状态} & \textbf{约束} \\ \hline
待接单 & 骑手接单 & 已接单 & 骑手已审核、在线且无超出上限的活动任务；并发时仅一人成功。 \\ \hline
已接单 & 确认到店 & 已接单 & 记录到店时间，不额外增加状态。 \\ \hline
已接单 & 确认取餐 & 已取餐 & 商家已出餐，操作人为当前绑定骑手。 \\ \hline
已取餐 & 开始配送 & 配送中 & 展示必要收货信息和地图导航入口。 \\ \hline
配送中 & 确认送达 & 已送达 & 记录送达时间，完成后限制骑手继续读取敏感地址。 \\ \hline
已接单/已取餐/配送中 & 上报异常 & 异常待处理 & 必须选择异常类型并填写说明，不允许骑手直接取消订单。 \\ \hline
异常待处理 & 恢复/改派/终止 & 原状态/已取消 & 由管理员记录处理结果；改派时终止原任务并建立新任务。 \\ \hline
\end{longtable}

\subsubsection{订单与配送联动}

\begin{longtable}{|L{0.20\linewidth}|L{0.22\linewidth}|L{0.20\linewidth}|L{0.28\linewidth}|}
\caption{订单与配送联动矩阵}\label{tab:order-delivery-map}\\ \hline
\textbf{订单状态} & \textbf{配送状态} & \textbf{可操作角色} & \textbf{主要操作} \\ \hline
待支付/待商家接单/制作中 & 无 & 顾客/商家 & 支付、取消、接单、拒单、出餐 \\ \hline
待骑手接单 & 待接单 & 骑手 & 查看任务、接单 \\ \hline
待取餐 & 已接单 & 骑手/商家 & 到店、核对、取餐 \\ \hline
配送中 & 已取餐/配送中 & 骑手 & 导航、送达、上报异常 \\ \hline
已送达/已完成 & 已送达 & 顾客 & 确认收货、评价 \\ \hline
已取消 & 无/已取消 & 系统/管理员 & 释放资源、保存取消和处理记录 \\ \hline
\end{longtable}

\subsubsection{结算与资产回退顺序}

\srsfigure{figures/srs/07-settlement-sequence.pdf}{订单结算与资产处理顺序图}{UML 顺序图}{fig:settlement-sequence}

结算顺序固定为“商品小计→会员折扣→优惠券→积分抵扣→钱包/模拟支付”。每单最多使用一张优惠券；优惠后应付金额不得小于零；订单取消时按原流水生成反向流水，不直接删除原记录。

\subsection{AI 客服与智能点餐流程}

\srsfigure{figures/srs/08-ai-assistant-sequence.pdf}{AI 客服与智能点餐顺序图}{UML 顺序图}{fig:ai-assistant}

AI 客服可回答平台规则，也可在用户登录且后端通过资源归属校验后查询当前用户的订单状态。智能点餐根据预算、口味、忌口和当前可售菜单返回结构化候选商品；候选结果必须经商品服务二次校验，并由顾客明确确认后才能加入购物车。AI 不得直接支付、取消订单、修改库存或调整用户资产。
